\chapter{Problèmes de synthèse}
Les synthèses S1 à S4 et S8 mobilisent la spécialité et ses approfondissements signalés ; S5 à S7 mobilisent aussi les mathématiques expertes.
\section{S1 -- Héron : convergence et efficacité}
Pour $a>0$, $u_0>0$, poser $u_{n+1}=\tfrac12(u_n+a/u_n)$.
\begin{enumerate}
\item Montrer que la suite est définie et que $u_n\ge\sqrt a$ pour $n\ge1$.
\item Établir la décroissance à partir du rang 1, puis la convergence vers $\sqrt a$.
\item Démontrer $u_{n+1}-\sqrt a=(u_n-\sqrt a)^2/(2u_n)$.
\item Pour $a=2,u_0=2$, donner un encadrement rationnel de $\sqrt2$ après trois itérations. Expliquer comment certifier une précision sans connaître la racine.
\item Relier la méthode à Newton et expliquer pourquoi l'exemple ne justifie pas la convergence de Newton pour toute fonction.
\end{enumerate}

\section{S2 -- Géométrie d’un tétraèdre}
Dans un repère orthonormé, $O=(0,0,0)$, $A=(2,0,0)$, $B=(0,3,0)$, $C=(0,0,6)$ ; les coordonnées sont en centimètres.
\begin{enumerate}
\item Déterminer une équation du plan $(ABC)$ et un vecteur normal.
\item Calculer le projeté orthogonal $H$ de $O$ sur ce plan, puis $OH$.
\item Calculer le volume de $OABC$ en prenant $OAB$ pour base, puis l'aire de $ABC$ en utilisant la hauteur $OH$.
\item Montrer que la sphère de diamètre $[OC]$ a équation $x^2+y^2+z^2-6z=0$.
\item Déterminer son intersection avec le plan $z=3$ ; expliquer pourquoi aucune figure seule ne prouve ces résultats.
\end{enumerate}

\section{S3 -- Logarithme, intégrale et sommes}
Pour $n\ge1$, poser $L_n=\frac1n\sum_{k=0}^{n-1}(1+k/n)^{-1}$ et $R_n=\frac1n\sum_{k=1}^{n}(1+k/n)^{-1}$.
\begin{enumerate}
\item Justifier $R_n\le\ln2\le L_n$ et calculer $L_n-R_n$.
\item En déduire un algorithme certifiant une erreur au plus $10^{-4}$ par le milieu de l'encadrement.
\item Poser $J_n=\int_0^1x^n/(1+x)\,dx$. Montrer $0\le J_n\le1/(n+1)$ et $J_n+J_{n+1}=1/(n+1)$.
\item En déduire une expression de $\ln2$ par une somme alternée avec une borne de reste.
\item Comparer approximation, encadrement prouvé et valeur affichée en décimal.
\end{enumerate}

\section{S4 -- Surbooking : modèle et décision}
Une salle possède 100 places. Pour une simulation pédagogique, chaque personne inscrite vient indépendamment avec probabilité $0{,}9$ ; on inscrit $n\ge100$ personnes.
\begin{enumerate}
\item Définir la loi du nombre $X_n$ de présents, sa moyenne et sa variance.
\item Écrire exactement $P(X_n>100)$ et proposer un algorithme cherchant le plus grand $n$ tel que ce risque soit au plus $0{,}05$.
\item Justifier que ce risque croît avec $n$ et tend vers $1$, afin de justifier l'arrêt de la recherche.
\item Pour $n=105$, donner une majoration par Tchebychev et dire si elle décide la contrainte de $5\,\%$.
\item Expliquer quelles hypothèses devraient être examinées avant une application réelle.
\end{enumerate}

\section{S5 -- Complexes, rotations et polygones}
Poser $\omega=e^{2i\pi/5}$ et $A_k$ d'affixe $\omega^k$, $0\le k\le4$.
\begin{enumerate}
\item Prouver que les cinq points sont distincts et forment un pentagone régulier.
\item Calculer leur somme et leur barycentre.
\item Montrer $|1-e^{it}|=2|\sin(t/2)|$ ; en déduire le côté et la diagonale du pentagone.
\item Poser $s=\omega+\omega^{-1}$. À partir de la somme des racines, montrer $s^2+s-1=0$ et déterminer $\cos(2\pi/5)$.
\item Écrire la matrice de la rotation envoyant $A_k$ sur $A_{k+1}$ et montrer que sa cinquième puissance est l'identité.
\end{enumerate}

\section{S6 -- RSA : preuve et limites}
Prendre deux nombres premiers distincts $p=5,q=11$, $N=pq$, $e=3$.
\begin{enumerate}
\item Trouver $d\in\{1,\ldots,39\}$ tel que $ed\equiv1\pmod{40}$.
\item Chiffrer $m=7$ par $c\equiv m^e\pmod N$, puis déchiffrer par $c^d\pmod N$ sans développer de très grands entiers.
\item Prouver que $m^{ed}\equiv m$ modulo $5$ puis modulo $11$ pour tout entier $m$, y compris s'il est divisible par l'un de ces nombres.
\item En déduire la correction modulo $55$ à l'aide de Gauss.
\item Pourquoi cet exemple ne fournit-il pas un système sécurisé prêt à l'emploi ?
\end{enumerate}

\section{S7 -- Trois états et convergence}
Une chaîne de Markov possède la matrice de transition
$P=\begin{pmatrix}1/2&1/4&1/4\\1/4&1/2&1/4\\1/4&1/4&1/2\end{pmatrix}$, avec des distributions lignes.
\begin{enumerate}
\item Vérifier que $P$ est une matrice de transition et que $\pi=(1/3,1/3,1/3)$ est invariante.
\item Noter $J$ la matrice dont tous les coefficients valent $1$ et $Q=J/3$. Montrer $Q^2=Q$ et $P=Q+\tfrac14(I-Q)$.
\item En déduire $P^n=Q+4^{-n}(I-Q)$ pour $n\ge0$.
\item Déterminer la distribution au rang $n$ pour une distribution initiale quelconque $\mu_0$ ; prouver convergence et unicité de l'invariante.
\item Trouver un rang garantissant que chaque coordonnée diffère de $1/3$ d'au plus $10^{-3}$, quel que soit $\mu_0$.
\end{enumerate}

\section{S8 -- Une fonction définie par une intégrale}
Pour $x\ge0$, poser $F(x)=\int_0^xe^{-t^2}\,dt$.
\begin{enumerate}
\item Justifier que $F$ est dérivable, strictement croissante, puis concave sur $[0,+\infty[$.
\item Pour $x\in[0,1]$, montrer $x-x^3/3\le F(x)\le x$.
\item Pour $x\ge1$, utiliser $t^2\ge t$ sur $[1,x]$ pour majorer $F(x)$ indépendamment de $x$. En déduire l'existence d'une limite finie $L$.
\item Montrer $0\le L-F(x)\le e^{-x}$ pour $x\ge1$.
\item Expliquer pourquoi cela ne donne pas une primitive élémentaire de $e^{-x^2}$, ni la valeur exacte de $L$. Relier la question à une future théorie des densités.
\end{enumerate}
